El presente Trabajo Fin de Máster describe el diseño y la implementación de una plataforma analítica reproducible para estudiar los patrones de movilidad urbana en la ciudad de Nueva York y predecir la demanda de servicios de taxi y vehículos de alquiler con conductor. La fuente principal son los registros de viajes publicados por la New York City Taxi and Limousine Commission (TLC), integrando las modalidades Yellow Taxi, Green Taxi, For-Hire Vehicle (FHV) y High Volume For-Hire Vehicle (HVFHV). El periodo de estudio comprende 2018--2025, de modo que el análisis permite comparar la movilidad previa a la pandemia de COVID-19, la disrupción observada durante 2020 y el proceso de recuperación posterior.

La plataforma se estructura siguiendo una arquitectura de medallón. La capa Bronze conserva los ficheros de origen en formato Parquet; las capas Silver y Gold se materializan en PostgreSQL y se transforman mediante dbt. La ingesta, planificación y control operacional del flujo se centralizan en Apache NiFi, que organiza el procesamiento mediante grupos de procesos, parámetros, rutas de éxito y error, reintentos y mecanismos de procedencia del dato. Docker Compose encapsula los servicios necesarios y permite desplegar de forma reproducible NiFi, PostgreSQL, dbt, el entorno Python de machine learning y Apache Superset. La capa Gold contiene dimensiones, tablas de hechos y \textit{data marts} diseñados para responder a preguntas de negocio relacionadas con demanda, movilidad, congestión aproximada, rendimiento económico y oportunidades operativas para conductores.

La predicción se formula como un problema de regresión sobre el número de viajes iniciados por zona TLC y hora. Se consideran horizontes de una hora y veinticuatro horas y se plantea la comparación de tres aproximaciones complementarias: regresión de Poisson, Random Forest y LightGBM\@. La evaluación utiliza particiones temporales para evitar fuga de información y métricas MAE, RMSE y WAPE\@. La memoria separa explícitamente los resultados observados de las estimaciones y simulaciones de rentabilidad, y reserva las cifras experimentales definitivas para la ejecución reproducible del repositorio final.
